%% SYNTHETIC TEST UNIT generated by writing/tools/make_outline_book.py from FINAL_OUTLINE_v2 -- not book text
\chapter{硬遍举过程}\label{ch:excl}
本单元为测试用的合成文本（由 FINAL\_OUTLINE\_v2 自动生成），陈述内容从略。

\section{第 1 节（测试）}\label{sec:excl-da}

\begin{definition}[介子分布振幅]\label{def:excl-da}
（测试用合成陈述：内容从略。）
\uses{def:prelim-wilson, def:prelim-gamma5}
\end{definition}

\begin{proposition}[单圈 ERBL 演化 \tagR]\label{prop:excl-erbl}
（$R_{\mathrm{PT}}$：在重整化微扰论中对红外正规化的部分子态严格。）
（测试用合成陈述：内容从略。）
非微扰地位：\cref{asm:intro-wightman}、\cref{asm:intro-composite}。
\uses{def:excl-da, prop:pdf-lo-kernels, asm:intro-wightman, asm:intro-composite}
\end{proposition}

\begin{proposition}[色单态介子的软胶子在被积函数层次相消 \tagR]\label{prop:excl-soft-cancel}
（$R_{\mathrm{cond}}$：在所列编号假设之下证明完整；结论的地位不高于其中最弱的假设。）
前提：\cref{asm:appN-endpoint}、\cref{asm:appN-collins}、\cref{asm:appN-ward-product}。
（测试用合成陈述：内容从略。）
\uses{def:excl-da, def:method-soft-approx, asm:appN-endpoint, asm:appN-collins, asm:appN-ward-product, def:appN-excl-approx}
\end{proposition}

\section{第 2 节（测试）}\label{sec:excl-soft}

\begin{theorem}[介子电磁形状因子的因子化 \tagA]\label{thm:excl-pion-ff}
（测试用合成陈述：内容从略。）
非微扰地位：\cref{asm:intro-bridge}。
\uses{def:excl-da, prop:excl-soft-cancel, thm:method-template, asm:appN-endpoint, asm:intro-bridge}
\end{theorem}
\begin{proofsketch}测试。\end{proofsketch}
\paragraph{地位与缺口}测试：本陈述的地位由\cref{asm:appN-endpoint}、\cref{asm:appN-collins}、\cref{asm:intro-bridge}、\cref{prop:excl-endpoint}决定。

\begin{proposition}[领头阶硬核 \tagR]\label{prop:excl-pion-lo}
（$R_{\mathrm{PT}}$：在重整化微扰论中对红外正规化的部分子态严格。）
（测试用合成陈述：内容从略。）
\uses{def:excl-da, def:prelim-feynman-rules, lem:appA-fierz}
\end{proposition}

\begin{example}[领头阶 $\pi$ 形状因子]\label{ex:excl-pion-ff-lo}
（测试用合成陈述：内容从略。）
\uses{prop:excl-pion-lo}
\end{example}

\section{第 3 节（测试）}\label{sec:excl-ff}

\begin{theorem}[$\gamma^*\gamma\to\pi^0$ 形状因子 \tagA]\label{thm:excl-transition}
（测试用合成陈述：内容从略。）
非微扰地位：\cref{asm:intro-bridge}。
\uses{def:excl-da, thm:method-template, asm:intro-bridge}
\end{theorem}
\begin{proofsketch}测试。\end{proofsketch}
\paragraph{地位与缺口}测试：见文献。

\begin{proposition}[端点区域与软重叠 \tagP]\label{prop:excl-endpoint}
（测试用合成陈述：内容从略。）
\uses{thm:excl-pion-ff}
\end{proposition}
\paragraph{地位与缺口}测试：见文献。

\begin{definition}[广义部分子分布]\label{def:excl-gpd}
（测试用合成陈述：内容从略。）
\uses{def:pdf-quark, def:pdf-lightray}
\end{definition}

\section{第 4 节（测试）}\label{sec:excl-endpoint}

\begin{proposition}[多项式性与 Ji 求和规则 \tagR]\label{prop:excl-polynomiality}
（$R_{\mathrm{def}}$：由算符定义直接推出。）
（测试用合成陈述：内容从略。）
非微扰地位：\cref{asm:intro-wightman}、\cref{asm:intro-composite}。
\uses{def:excl-gpd, asm:intro-wightman, asm:intro-composite}
\end{proposition}

\begin{theorem}[DVCS 的因子化 \tagA]\label{thm:excl-dvcs}
（测试用合成陈述：内容从略。）
非微扰地位：\cref{asm:intro-bridge}。
\uses{def:excl-gpd, thm:method-template, asm:intro-bridge}
\end{theorem}
\begin{proofsketch}测试。\end{proofsketch}
\paragraph{地位与缺口}测试：见文献。

\begin{proposition}[DVCS 领头阶系数函数 \tagR]\label{prop:excl-dvcs-lo}
（$R_{\mathrm{PT}}$：在重整化微扰论中对红外正规化的部分子态严格。）
（测试用合成陈述：内容从略。）
\uses{def:excl-gpd, def:dis-compton, def:prelim-feynman-rules}
\end{proposition}

\section{第 5 节（测试）}\label{sec:excl-gpd}

\begin{theorem}[纵向光子介子产生的因子化 \tagA]\label{thm:excl-meson-L}
（测试用合成陈述：内容从略。）
非微扰地位：\cref{asm:intro-bridge}。
\uses{def:excl-da, def:excl-gpd, prop:excl-soft-cancel, thm:method-template, asm:appN-endpoint, asm:appN-collins, asm:intro-bridge}
\end{theorem}
\begin{proofsketch}测试。\end{proofsketch}
\paragraph{地位与缺口}测试：本陈述的地位由\cref{asm:appN-endpoint}、\cref{asm:appN-collins}、\cref{asm:intro-bridge}决定。

\begin{proposition}[横向光子：标准共线因子化失效 \tagV]\label{prop:excl-meson-T}
（测试用合成陈述：内容从略。）
\uses{thm:excl-meson-L}
\end{proposition}
\paragraph{地位与缺口}测试：见文献。

\begin{remark}[宽角强子散射]\label{rem:excl-hadron-scattering}
（测试用合成陈述：内容从略。）
% uses: none
\end{remark}

\section{第 6 节（测试）}\label{sec:excl-dvcs}

\begin{proposition}[近阈重夸克偶素光生（重夸克极限） \tagP]\label{prop:excl-near-threshold}
（测试用合成陈述：内容从略。）
\uses{def:excl-gpd}
\end{proposition}
\paragraph{地位与缺口}测试：见文献。
